\begin{table*}[t!]
  \centering
  \caption{\textbf{Leave-subunit-out recovery of $175$ obligate heteromeric complex subunits in six prokaryotic genomes under two decoy designs.} In each trial the raw score of the true subunit's claim for the complex term is set to $0.50$ (two-subunit complexes) or $0.45$ (larger complexes), its partner subunits are raised to raw scores $\ge 0.85$, and five decoy proteins without annotation to the complex term receive synthetic domain-channel claims drawn from $U(0.42, 0.49)$. In the \emph{unrelated-loci} design the decoys are random proteins from loci unrelated to the complex; in the \emph{within-operon} design the same decoy scores are given to proteins from the true subunit's own operon, completed from its co-directional gene run and, where still fewer than five candidates were available (71 of 175 trials), from the nearest genes on the same contig. Real-channel rows rank the candidates by their real DIAMOND or InterProScan/Pfam claims for the complex term, capped at the trial's attenuation level. Hit@1 scores ties at the top as $1/\text{ties}$; precision, recall, and $F_1$ are computed at threshold $0.50$ among the six candidates, and complex coherence is genome-wide. Stage 3 is the guarded tool-calling agent (Qwen3.8-27B served by vLLM, temperature $0$, single run with seed $42$). Stage 2A alone changes no score on this benchmark (its output is identical to coherence + operon priors) and is not shown. Rows are given for all six species ($n = 175$), the five non-calibration species ($n = 145$), and \textit{Synechocystis} sp.\ PCC 6803 alone ($n = 25$); bold marks the column maximum within each design and cohort. Exact McNemar tests on the strict Hit@1 indicator are given in the Supplementary Information.}
  \label{tab:complex_lso}
  \resizebox{\textwidth}{!}{%
  \begin{tabular}{llcccccc}
    \toprule
    \textbf{Cohort} & \textbf{Baseline or Method Configuration} & \textbf{Subunit Trials ($n$)} & \textbf{Hit@1} & \textbf{Precision@0.5} & \textbf{Recall@0.5} & \textbf{$F_1$@0.5} & \textbf{Complex Coherence} \\
    \midrule
